\documentclass[11pt,a4paper]{article}
\usepackage[T1]{fontenc}
\usepackage[utf8]{inputenc}
\usepackage[main=italian,provide=*]{babel}
\usepackage{amsmath,amssymb}
\usepackage{geometry}
\geometry{margin=2.3cm,top=2.5cm,bottom=2.7cm}
\usepackage{booktabs}
\usepackage[table,dvipsnames]{xcolor}
\usepackage{tcolorbox}
\tcbuselibrary{skins,breakable}
\usepackage{titlesec}
\usepackage{enumitem}
\usepackage{needspace}
\usepackage{fancyhdr}
\usepackage{graphicx}
\usepackage{hyperref}
\hypersetup{colorlinks=true,linkcolor=Sepia,citecolor=Sepia,urlcolor=Sepia}
\newcommand{\Tr}{\operatorname{Tr}}
\newcommand{\ket}[1]{\lvert #1 \rangle}
\newcommand{\bra}[1]{\langle #1 \rvert}

\definecolor{inkblue}{HTML}{1B3A5C}
\definecolor{lightblue}{HTML}{EAF1F8}
\definecolor{accent}{HTML}{B0562A}
\definecolor{softgray}{HTML}{6B6B6B}
\definecolor{okgreen}{HTML}{2E6B3E}
\definecolor{okgreenbg}{HTML}{EAF5EC}

\titleformat{\section}
  {\normalfont\Large\bfseries\color{inkblue}}
  {\thesection}{0.6em}{}[{\color{inkblue!40}\titlerule}]
\titlespacing{\section}{0pt}{0.9em}{0.4em}
\titleformat{\subsection}
  {\normalfont\large\bfseries\color{accent}}
  {}{0em}{}
\titlespacing{\subsection}{0pt}{0.6em}{0.2em}

\pagestyle{fancy}
\fancyhf{}
\renewcommand{\headrulewidth}{0.4pt}
\fancyhead[L]{\small\color{softgray}VQE con termine DM sotto rumore}
\fancyhead[R]{\small\color{softgray}\thepage}
\fancyfoot[C]{\small\color{softgray}Tesi triennale, Samuele --- Relatore: Prof.\ Alessandro Chiesa}

\newtcolorbox{keyeq}[1][]{
  colback=lightblue, colframe=inkblue!70, boxrule=0.6pt,
  arc=2pt, left=8pt, right=8pt, top=3pt, bottom=3pt, #1
}
\newtcolorbox{okbox}[1]{
  colback=okgreenbg, colframe=okgreen!60, boxrule=0.6pt,
  arc=2pt, left=7pt, right=7pt, top=2pt, bottom=2pt,
  fonttitle=\bfseries\color{okgreen}, title={#1}
}

\title{\vspace{-1.5em}\color{inkblue}VQE con termine DM sotto rumore\\[0.2em]
\Large\normalfont\color{softgray}Quanto degrada un risultato ideale su un simulatore rumoroso}
\author{Note di lavoro --- Tesi triennale, Samuele\\ Relatore: Prof.\ Alessandro Chiesa}
\date{}

\begin{document}
\sloppy
\maketitle
\thispagestyle{fancy}
\vspace{-2.0em}

File di codice: \texttt{vqe\_dm\_rumoroso\_dimero.py}. Punto di lavoro:
$b/J=0.35$, $D/J=0.80$, $J=1$ (``test 2''). Modello di rumore: il
canale depolarizzante 1q/2q + readout simmetrico di
\texttt{noise\_model\_dimero.py}, ai valori di riferimento
\texttt{ibm\_torino} salvo dove diversamente indicato.

\tableofcontents

\section{Obiettivo}

Il VQE ideale (Parte 1, nessun rumore) trova parametri
$\theta^*_\text{ideale}$ che preparano lo stato fondamentale
dell'Hamiltoniana del dimero con DM entro la precisione numerica. Questo
documento chiede: se lo stesso circuito, con la stessa ricetta di
parametri, viene eseguito su un simulatore \textbf{rumoroso}, quanto
degradano energia e fedeltà? È la prima domanda della Parte 2 (rumore),
e la base su cui si costruiscono tutti gli stadi successivi.

\begin{keyeq}[title={Scelta di modellazione dichiarata}]
\centering
Si riusano i parametri VQE già ottimizzati nel caso ideale, \textbf{senza}
rioptimizzare sotto rumore: questo documento misura la degradazione di un
risultato ideale, non la convergenza di un VQE che veda il rumore durante
l'ottimizzazione (quella domanda è affrontata separatamente,
\texttt{risultati\_vqe\_noise\_aware\_definitivo.tex}).
\end{keyeq}

\section{Osservabili: energia e fedeltà rumorose}

Con il circuito rumoroso, lo stato preparato non è più puro: è descritto
da un operatore densità $\rho$. L'energia si generalizza in modo diretto,
$E_\text{rumoroso}=\Tr[\rho H]$. La fedeltà rispetto allo stato di
riferimento puro $\ket{\psi_\text{ideale}}$ (fedeltà di Uhlmann,
semplificata per un riferimento puro) è
\begin{keyeq}
\centering
$F(\rho,\ket{\psi_\text{ideale}}) = \bra{\psi_\text{ideale}}\rho\ket{\psi_\text{ideale}}$
\end{keyeq}
--- verificata contro \texttt{qiskit.quantum\_info.state\_fidelity} su
uno stato misto casuale (non il caso specifico di questo problema, per
evitare una verifica circolare).

\section{Risultati}

\begin{center}
\renewcommand{\arraystretch}{1.3}
\begin{tabular}{@{}lccc@{}}
\toprule
\textbf{Condizione} & $\boldsymbol{E}$ & $\boldsymbol{F}$ & \textbf{CNOT}\\
\midrule
Rumore nullo (verifica di coerenza) & $-3.57321145$ & $1.0000000000$ & $2$\\
Rumore di riferimento (\texttt{ibm\_torino}) & $-3.50164605$ & $0.9850060074$ & $2$\\
\bottomrule
\end{tabular}
\end{center}

\begin{okbox}{Verificato}
Il caso a rumore nullo riproduce $E_0=-3.57321145$ (l'energia esatta,
Parte 1) a precisione numerica --- coerenza interna della pipeline. Al
rumore di riferimento: scostamento di energia $+0.071565$, perdita di
fedeltà $1-F=0.014994$ ($\sim1.5\%$) --- con un ansatz a soli 2 CNOT,
poco margine per il rumore di accumularsi.
\end{okbox}

\section{Dipendenza da \texorpdfstring{$\varepsilon_{2q}$}{eps\_2q}}

\begin{figure}[h]
\centering
\includegraphics[width=0.92\linewidth]{figure/fig_passo2_scan_eps2q.pdf}
\caption{\textit{Energia e fedeltà rumorose in funzione di
$\varepsilon_{2q}$ (a $\varepsilon_{1q}=2.9\times10^{-4}$,
$p_\text{readout}=2.3\times10^{-2}$ fissati). La linea tratteggiata
segna il valore di riferimento $\varepsilon_{2q}=3.8\times10^{-3}$.
Degradazione monotona e liscia, nessuna soglia o comportamento anomalo
nell'intervallo esplorato.}}
\end{figure}

\section{Cosa non è ancora stato controllato qui}

Questo documento fissa un solo punto (o uno scan di un solo parametro):
non dice ancora \emph{come} la fedeltà si comporti se il circuito viene
allungato (nessuna dinamica qui, solo lo stato fondamentale statico), né
se il rumore possa essere assorbito da un'ottimizzazione diversa dei
parametri (affrontato separatamente, VQE noise-aware). Serve solo come
riferimento di degradazione per uno stato preparato, prima di introdurre
qualunque evoluzione temporale.

\section{Considerazioni e cosa serve al passo successivo}

Due fatti da questo documento sono usati esplicitamente nel passo
successivo (quantum simulation Trotter sotto rumore): (1) i parametri
$\theta^*_\text{ideale}$ e il loro CNOT-count minimo (2, per la sola
preparazione) sono il punto di partenza --- il passo successivo aggiunge
$N$ passi di Trotter sopra questa stessa preparazione; (2) la scala della
degradazione qui osservata ($1-F\sim1.5\%$ al rumore di riferimento) è
il termine di \emph{preparazione} che il passo successivo isolerà come
contributo costante, indipendente da $N$, nella fedeltà totale --- un
fatto verificato esplicitamente lì, non solo assunto qui.

\begin{okbox}{Sintesi in una riga}
Riusare parametri VQE ideali sotto rumore di riferimento costa
$\sim1.5\%$ di fedeltà con un ansatz a 2 CNOT --- la scala di riferimento
per tutti gli stadi a valle.
\end{okbox}

\needspace{4\baselineskip}
\section*{File prodotti}
\addcontentsline{toc}{section}{File prodotti}

\texttt{vqe\_dm\_rumoroso\_dimero.py}, questo documento.

\end{document}
